%%%%%%%%%%%%%%%%%%%%%%%%%%%%%%%%%%%%%%%%%%%%%%%%%%%%%%%%%%%%%%%%%%%%%%%%%%%%%%%%
% BLOQUE 7 — VÍNCULO CON LA PROPUESTA TÉCNICO-ECONÓMICA
% El enunciado deja a criterio del grupo si esto va dentro de
% Desarrollo o de Conclusiones; se dejó como sección propia (entre
% Entregables y Tendencias) para que el orden de la pauta sea
% explícito. Si el equipo prefiere fusionarlo con Desarrollo o
% Conclusiones, basta con mover este \input en main.tex y ajustar el
% \section por \subsection.
%%%%%%%%%%%%%%%%%%%%%%%%%%%%%%%%%%%%%%%%%%%%%%%%%%%%%%%%%%%%%%%%%%%%%%%%%%%%%%%%

\section{Vínculo con la propuesta técnico-económica}

% TODO (equipo): este es "el caso más claro de arquitectura que cambia
% la estructura de costos de CAPEX a OPEX puro" (ficha TI-05). Discutir
% ese cambio y el riesgo de "la factura que crece con el éxito del
% sistema" (a mayor tráfico/adopción, mayor costo variable, a
% diferencia de una inversión fija ya amortizada). Conectar con el
% modelo de costos del bloque 6.

\clearpage
